\begin{activity}\label{A:0.5.2}
Suppose the following sinusoidal function models the water level on a pier in the ocean as
it changes due to the tides during a certain day.
\[ w(t) = 4.3 \sin \left( 0.51 t + 0.82 \right) + 10.6 \]
\ba
\item Using the formula above, make a table showing the water level every two hours for a 24 hour period starting at midnight.
    \begin{center}
        \begin{tabular}[h!]{|c||c|c|c|c|c|c|c|c|c|c|c|c|c|}
            \hline
            time (hours) & 0 & 2 & 4 & 6 & 8 & 10 & 12 & 14 & 16 & 18 & 20 & 22 & 24 \\
            \hline
            water level (ft) & & & & & & & & & & & & & \\ \hline
        \end{tabular}
    \end{center}
\item Sketch a graph of this function using the data from your table in part (a).
    \begin{center}
%         \begin{tikzpicture}
%             \begin{axis}[ymin=0, ymax=20, xmin=0, xmax=24, domain=0:24,
%                 grid, xtick={2,4,6,8,10,12,14,16,18,20,22,24},
%             ytick={0,2,4,6,8,10,12,14,16,18,20}, xlabel={time (hours)}, ylabel={water level
%             (ft)}, xscale=1.25]
%                 \addplot[smooth] {0*x};
%             \end{axis}
%         \end{tikzpicture}
        \includegraphics[width=0.6\columnwidth]{figures/0-5-fig10.pdf}
    \end{center}
\item What is the period of oscillation of this function? % Period $ = 2 \pi / B \approx  12.3$ hours.
\item What time is high tide?  % The sine function reaches a peak when its input is $\pi/2$.  So $0.51 t + 0.82 = \pi/2$ and $t = 1.47$ hours, or about 1:30 am.

\ea
\end{activity}\aftera
